On the domain $\Omega=\{(x,y)\in\mathbb{R}^2 \mid 0 < x < \frac{\pi}2\}$,
solve the
partial differential equation
\begin{equation}
\frac{\partial u}{\partial x}
+
u\frac{\partial u}{\partial y}
=
-y
\end{equation}
with the boundary condition
\[
u(0,y_0) = f(y_0)\quad \text{for $y_0\in\mathbb{R}$}.
\]
\begin{teilaufgaben}
\item
Is the problem well posed?
\item
Find an explicit solution in the form $u(x,y)$ for $f(y_0)=y_0$.
\end{teilaufgaben}

\begin{loesung}
This is a quasilinear partial differential equation of first order which
can be solved using the method of characteristics.
The system of ordinary differential equations for the characteristics
is
\begin{align*}
\dot{x}(t) &=  1 \\
\dot{y}(t) &=  u(t) \\
\dot{u}(t) &= -y(t).
\end{align*}
The first equation has the solution $x(t)=t+x_0$, the boundary condition
implies that $x_0=0$.
The variables $t$ and $x$ are thus interchangable.

Combining the last two equations gives the coupled second order differential
equations
\begin{align*}
\ddot{y} &= \phantom{-}\dot{u} = -y \\
\ddot{u} &=          - \dot{y} = -u,
\end{align*}
both of which have linear combinations of the functions $\sin t$ and
$\cos t$.
Thus we set $y(t) = A\cos t + B \sin t$ and derive from the second
equation, that $u(t) = \dot{y}(t) = -A\sin t + B \cos t$.

The boundary condition at $x=0$ implies $y_0=y(0)=A$ and thus
$f(y_0) = u(0) = B$.
Thus the solutions with $t$ already eliminated are
\begin{align*}
y(x) &= f(y_0)\cos x +   y_0  \sin x \\
u(x) &=   y_0 \cos x - f(y_0) \sin x 
\end{align*}
For $f(y_0)=y_0$ these become
\begin{align}
y(x) &=   y_0(\cos x + \sin x) 
&&\Rightarrow &
y_0 &= \frac{y}{\cos x + \sin x}
\label{30000033:eqn:ysol}
\\
u(x) &=   y_0(\cos x - \sin x)
&&\Rightarrow &
u &= y\frac{\cos x - \sin x}{\cos x + \sin x},
\label{30000033:eqn:usol}
\end{align}
giving the solution formula requested in b).
\begin{figure}
\centering
\includeagraphics[]{graph.pdf}
\caption{Projection of characteristic curves into the domain of the
differential equation of problem \ref{30000033}.
\label{30000033:char}}
\end{figure}
The questions asked can now be answered directly from the solution formulae
\eqref{30000033:eqn:ysol} and \eqref{30000033:eqn:usol}.
\begin{teilaufgaben}
\item
Since the denominator in \eqref{30000033:eqn:ysol} does not vanish
in the interval $(0,\frac{\pi}2)$, there is a well defined characteristic
curve for each point $(x,y)$ that reaches the point.
Also all the characteristic curves are obtained by stretching along
the $y$-direction, no point is reached via two different characteristics.
Therefore, the problem is well posed.
\item
The solution is given by \eqref{30000033:eqn:usol}.
Using $\cos x = \sin(\frac{\pi}2-x)$ and the identity
\[
\sin\alpha\pm\sin\beta
=
2\sin\frac{\alpha\pm\beta}2 \cos\frac{\alpha\mp\beta}2
,
\]
we can further simplify
\begin{align*}
\cos x \pm \sin x
&=
\sin\biggl(\frac{\pi}2-x\biggr) \pm \sin x
=
2
\sin\bigl(\frac12\biggl(\frac{\pi}2-x \pm x\biggr)\biggr)
\cos\bigl(\frac12\biggl(\frac{\pi}2-x \mp x\biggr)\biggr)
\\
&=
\begin{cases}
2
\sin \frac{\pi}4
\cos \bigl(\frac{\pi}4-x\bigr)
=
\sqrt{2}\cos\bigl(\frac{\pi}4-x\bigr)
\\
2
\cos \frac{\pi}4
\sin \bigl(\frac{\pi}4-x\bigr)
=
\sqrt{2}\sin\bigl(\frac{\pi}4-x\bigr)
\end{cases}
\qedhere
\end{align*}
\end{teilaufgaben}
\end{loesung}

\begin{diskussion}
Note that figure~\ref{30000033:char} also shows how the characteristics
no longer cover the domain if it is extended in the $x$-direction beyond
$x=\frac{3\pi}{4}$.
At this point, the problem would no longer be well posed, as the values
of the function $u(x,y)$ in the points $(\frac{3\pi}4,y)$ cannot be
determined from boundary conditions.
Not surprisingly, this coincides with the fact that at this point, the
denominator
\[
\cos\frac{3\pi}{4} + \sin\frac{3\pi}{4} 
=
-
\frac{\!\sqrt{2}}2
+
\frac{\!\sqrt{2}}2
=
0
\]
of the solution formula \eqref{30000033:eqn:usol} vanishes.
\end{diskussion}

\begin{bewertung}
Differential equations for characteristics ({\bf D}) 1 point,
solution for $x(t)$ including $x_0=0$ ({\bf X}) 1 point,
solutions for $y(t)$ und $u(t)$ ({\bf YU}) 1 point,
using boundary condition $y(0)=1$ to determine $A=1$ ({\bf A}) 1 point,
using boundary condition $u(0)=y_0$ to determine $B=y_0$ ({\bf B}) 1 point,
eliminate $y_0$ and $t$ giving $u(x,y)$ ({\bf U}) 1 point.
A common error was losing the minus sign in the third component of the
differential equation.
This leads to linear combinations of $\cosh x$ and $\sinh x$ instead
of $\cos x$ and $\sin x$.
However, determining the coefficients $A$ and $B$ and elimination of
$y_0$ still work the same, so partial credit was given for such solutions.
\end{bewertung}



